\section*{الفصل \ref{ch:b1:genomes} --- مورثومات الخلايا والفيروسات}

\begin{solution}{exo:b1:genomes:1}
البكتيريا: \omterm{def:b1:genomes:genome}{صبغي} حلقي واحد (مع \omterm{def:b1:genomes:genome}{بلازميدات}) في \omterm{def:b1:genomes:genome}{نويد} بلا غشاء، مضغوط
\omterm{prop:b1:genomes:supercoiling}{بالالتفاف الفائق} وببروتينات أساسية صغيرة، ونحو \qty{90}{\%} منه
مشفِّر. \omterm{def:b1:cell-unit-of-life:prokeuk}{وحقيقية النواة}: عدة \omterm{def:b1:genomes:genome}{صبغيات} خطية في نواة، ملفوفة على هستونات في
\omterm{def:b1:genomes:nucleosome}{صبغين}، وبضعة في المئة منها مشفِّرة.
\end{solution}

\begin{solution}{exo:b1:genomes:2}
147 \omterm{thm:b1:nucleic-acids:helix}{زوج قواعد} ملفوفة 1.7 مرة حول ثماني هستونات (اثنتان من كل من H2A
وH2B وH3 وH4)، تختمها H1، واحدة كل 200 زوج. والخرزات على الخيط
$\times 6$؛ والليف \qty{30}{nm} $\times 40$؛ والعرى $\times 1000$؛
\omterm{def:b1:genomes:genome}{وصبغي} الطور الاستوائي $\times \num{10000}$.
\end{solution}

\begin{solution}{exo:b1:genomes:3}
\qty{1.5}{\%}؛ والعناصر القافزة وبقاياها، \qty{45}{\%}.
\end{solution}

\begin{solution}{exo:b1:genomes:4}
\omterm{def:b1:genomes:genome}{مورثوم} (\omterm{def:b1:nucleic-acids:chain}{DNA} أو \omterm{def:b1:nucleic-acids:chain}{RNA}) في \omterm{def:b1:genomes:virus}{قفيصة} بروتينية، مغلَّف أحيانًا، لا يتكاثر إلا
داخل \omterm{def:b1:cell-unit-of-life:cell}{خلية} بآلة \omterm{def:b1:cell-unit-of-life:cell}{الخلية}. وخارج \omterm{def:b1:cell-unit-of-life:cell}{خلية} لا استقلاب له ولا نمو ولا استجابة
--- بل جسيم خامل؛ وهو «حي» بمعنى حمله معلومة وراثية تتطور فقط.
\end{solution}

\begin{solution}{exo:b1:genomes:5}
$\num{4.6e6}\times 0.34 = \qty{1.56}{mm}$؛ و$\num{4.6e6}/200 = \num{23000}$
\omterm{def:b1:genomes:nucleosome}{جسيم نووي}. وبدل ذلك: التفاف فائق سالب بالجيراز وطي في نحو خمسين عروة
على لب بروتيني، مع \omterm{def:b1:proteins:peptide}{بروتينات} أساسية صغيرة تثني \omterm{def:b1:nucleic-acids:chain}{DNA} --- أي ألف ضعف، في
\omterm{def:b1:genomes:genome}{نويد} بميكرومتر.
\end{solution}

\begin{solution}{exo:b1:genomes:6}
\omterm{def:b1:genomes:gene}{الإكسونات} $1500/8 = \qty{190}{bp}$؛ \omterm{def:b1:genomes:gene}{والإنترونات}
$25\,500/7 = \qty{3640}{bp}$؛ و$25.5/27 = \qty{94}{\%}$ من النسخة
مطروح.
\end{solution}

\begin{solution}{exo:b1:genomes:7}
\emph{E. coli} 930 \omterm{def:b1:genomes:gene}{مورّثة} لكل ميغا زوج؛ والخميرة 500؛ والذبابة 100؛
والإنسان 6؛ وسمكة الرئة 0.15. فكثافة \omterm{def:b1:genomes:gene}{المورّثات} تهبط أربع مراتب قدر
بينما يرتفع عددها خمسة أضعاف: فجزيء \omterm{def:b1:nucleic-acids:chain}{DNA} الزائد غير مورّثي.
\end{solution}

\begin{solution}{exo:b1:genomes:8}
$10^6\times 300 = \qty{3e8}{bp}$: أي \qty{9}{\%} من \omterm{def:b1:genomes:genome}{المورثوم}.
و$60\times
10^6/20 = \num{3e6}$ جيل: أي إدخال جديد كل ثلاثة أجيال،
بالمتوسط، في مكان ما من النسب.
\end{solution}

\begin{solution}{exo:b1:genomes:9}
البيضة تسهم \omterm{def:b1:cell-unit-of-life:cell}{بالسيتوبلازم} وبمتقدراته؛ \omterm{def:b1:eukaryotic-cell:mitochondrion}{ومتقدرات} النطفة القليلة تُستبعد
أو تُهدم بعد الإلقاح. فينتقل \omterm{def:b1:nucleic-acids:chain}{DNA} المتقدري من الأم إلى الولد بلا خلط،
وتتبع طفراته المتراكمة الأنساب الأمومية عبر الزمن.
\end{solution}

\begin{solution}{exo:b1:genomes:10}
في مزرعة مغذاة، على مدى ساعتين: يعطي الاندماج $2^6 = 64$ نسخة (واحدة
لكل انقسام مضيف)؛ ويعطي الحل 100 عاثية في \qty{40}{min}، وذريتها
$100^3 = 10^6$ في ساعتين إن كثر المضيفون --- فيفوز الحل بفارق كبير.
وفي مزرعة جائعة: يعطي الاندماج نسخة واحدة، آمنة داخل \omterm{def:b1:cell-unit-of-life:cell}{خلية} حية؛ ويعطي
الحل 100 عاثية بلا \omterm{def:b1:cell-unit-of-life:cell}{خلية} تعديها، فتضمحل. والمفتاح يقرأ حال المضيف لأن
أفضل استراتيجية تتوقف على هل ستتوفر مضائف جديدة.
\end{solution}

\begin{solution}{exo:b1:genomes:11}
$\num{1.3e11}\times 0.34\,\mathrm{nm} = \qty{44}{m}$ (أحادي الصيغة)؛
و$\num{6.5e8}$ \omterm{def:b1:genomes:nucleosome}{جسيم نووي}؛ والهستونات $\num{6.5e8}\times 108\,000\times
1.66\times 10^{-24} = \qty{117}{pg}$. والمناشئ: $\num{1.3e6}$، وكل
منها يضاعف \qty{50}{kb} على كل جانب بسرعة \qty{50}{bp/s}: أي \qty{1000}{s}
--- فالزمن ليس المشكلة، ما دامت المناشئ كافية؛ والكلفة هي
\omterm{def:b1:nucleic-acids:nucleotide}{النوكليوتيدات} (\qty{130}{pg} من \omterm{def:b1:nucleic-acids:chain}{DNA} لكل انقسام)، والهستونات، وزمن
صنعها، ونواة \omterm{def:b1:cell-unit-of-life:cell}{وخلية} أكبر بمرات، ولهذا تنقسم \omterm{def:b1:cell-unit-of-life:cell}{خلايا} كهذه ببطء.
\end{solution}

\begin{solution}{exo:b1:genomes:12}
التصميم يحوي ما يلزم ولا يزيد؛ والسجل يحوي ما جرى. \omterm{def:b1:genomes:gene}{فالمورّثات} الزائفة
جثث \omterm{def:b1:genomes:gene}{مورّثات} مضاعفة ماتت بالطفرة؛ والعناصر القافزة طفيليات نسخت نفسها
عشرات الملايين من السنين ولم تُزل قط؛ وعائلات \omterm{def:b1:genomes:gene}{المورّثات} تُظهر تضاعفًا
تلاه تباعد؛ \omterm{prop:b1:genomes:cvalue}{ومفارقة القيمة C} تُظهر أن مقدار \omterm{def:b1:nucleic-acids:chain}{DNA} يعكس تاريخ نسب من
التراكم والفقد لا حاجاته. والانتقاء يبقي \omterm{def:b1:genomes:gene}{المورّثات} عاملة ويترك الباقي
ينجرف: فما نقرؤه في \omterm{def:b1:genomes:genome}{مورثوم} تاريخ في معظمه.
\end{solution}

\begin{solution}{pb:b1:genomes:1}
\textbf{1.} $\num{6.4e9}\times 0.34\,\mathrm{nm} = \qty{2.18}{m}$.
\textbf{2.} $\num{6.4e9}/10 = \num{6.4e8}$ دورة.
\textbf{3.} $\num{2.5e8}\times 0.34 = \qty{8.5}{cm}$؛ و\qty{1.7}{cm}.
\textbf{4.} $\frac{4}{3}\pi\times 3^3 = \qty{113}{\micro m^3}$.
\textbf{5.} $\pi\times(1\,\mathrm{nm})^2\times 2.18\,\mathrm{m} =
\qty{6.8e-18}{m^3} = \qty{6.8}{\micro m^3}$: أي \qty{6}{\%} من \omterm{def:b1:eukaryotic-cell:nucleus}{النواة}.
\textbf{6.} $\num{6.4e9}\times 650\times 1.66\times 10^{-24} =
\qty{6.9e-12}{g} = \qty{6.9}{pg}$.
\textbf{7.} كتلة \omterm{def:b1:eukaryotic-cell:nucleus}{النواة} $113\times 10^{-12}\,\mathrm{mL}\times 1.1 =
\qty{124}{pg}$؛ \omterm{def:b1:nucleic-acids:chain}{وDNA} \qty{5.6}{\%} منها، \omterm{def:b1:proteins:peptide}{والبروتين} \qty{20}{\%}.
\textbf{8.} $\num{6.4e9}/200 = \num{3.2e7}$.
\textbf{9.} $\num{3.2e7}\times 108\,000\times 1.66\times 10^{-24} =
\qty{5.7}{pg}$: أي نحو كتلة \omterm{def:b1:nucleic-acids:chain}{DNA} نفسها.
\textbf{10.} $68/11 = 6.2$.
\textbf{11.} ست جسيمات نووية، أي \qty{1200}{bp} $= \qty{408}{nm}$ من
اللولب، في \qty{11}{nm}: أي بعامل $37$.
\textbf{12.} $2.18/37 = \qty{5.9}{cm}$ من الليف --- أي عشرة آلاف ضعف
قطر \omterm{def:b1:eukaryotic-cell:nucleus}{النواة}، مطويًا داخلها.
\textbf{13.} كلها، أي $\num{3.2e7}$، مع مثلها من الجديدة على النسخة
الثانية: أي $\num{6.4e7}$ في \qty{28800}{s}، أي نحو \num{2200} في
الثانية.
\textbf{14.} $75\,000\times 0.34/37 = \qty{690}{nm}$ من الليف لكل
عروة؛ و$\num{6.4e9}/75\,000 = \num{85000}$ عروة.
\textbf{15.} $\qty{8.5}{cm}/\qty{10}{\micro m} = 8500$.
\textbf{16.} $\pi\times 0.35^2\times 10 = \qty{3.85}{\micro m^3}$؛
\omterm{def:b1:nucleic-acids:chain}{وDNA} $\pi\times 10^{-6}\times 0.085\,\mathrm{m} = \qty{0.27}{\micro m^3}$:
أي \qty{7}{\%} --- والباقي \omterm{def:b1:proteins:peptide}{بروتين} هستوني وسقالي وماء.
\textbf{17.} مجموع \omterm{def:b1:nucleic-acids:chain}{DNA} $\num{6.4e9}$ زوج بالكثافة نفسها
($\num{2.5e8}$ زوج في \qty{3.85}{\micro m^3}، لكل صبيغية؛ وصبيغيتان
لكل \omterm{def:b1:genomes:genome}{صبغي}): $2\times 6.4\times 10^9/2.5\times 10^8\times
3.85 = \qty{197}{\micro m^3}$، وهو أكبر من نواة الطور البيني لأن
\omterm{def:b1:genomes:genome}{الصبغيات} صارت قضبانًا بينها فراغ.
\textbf{18.} \omterm{def:b1:genomes:nucleosome}{صبغين} الطور البيني لا بد أن يُقرأ ويُنسخ: فالبوليميرازات
وعواملها تحتاج إلى الوصول إلى \omterm{def:b1:nucleic-acids:chain}{DNA}، فيبقى معظمه في هيئة الليف
\qty{30}{nm} المفتوح والعرى؛ أما \omterm{def:b1:genomes:nucleosome}{صبغين} الطور الاستوائي فخامل ويمكن
رصه للنقل.
\textbf{19.} $\num{6.4e9}$ موقع عند $10^3$ في الثانية:
$\num{6.4e6}$ ثانية، أي 74 يومًا. \omterm{def:b1:proteins:peptide}{فالبروتينات} لا تبحث عشوائيًا: بل
ترتبط بجزيء \omterm{def:b1:nucleic-acids:chain}{DNA} ارتباطًا غير نوعي وتنزلق عليه، وتبحث نسخ كثيرة على
التوازي، فيُوجد الموقع في ثوان.
\textbf{20.} $\num{20000}\times 27\,\mathrm{kb} = \qty{540}{Mb}$، أي
\qty{17}{\%} من \qty{3.2}{Gb}؛ والمشفِّر $\num{20000}\times 1.5 =
\qty{30}{Mb}$، أي \qty{0.9}{\%} (ونحو 1.5\% بحساب \omterm{def:b1:genomes:gene}{المورّثات} الأقصر
حسابًا أدق).
\textbf{21.} $113\times 2\times 250/6400 = \qty{8.8}{\micro m^3}$
(نسختان)؛ و$\num{5e8}/8.8 = \qty{5.7e7}{bp/\micro m^3}$، أي كثافة
\omterm{def:b1:eukaryotic-cell:nucleus}{النواة} كلها نفسها.
\textbf{22.} $\num{2.6e11}/\num{5.7e7} = \qty{4600}{\micro m^3}$:
أي قطر \qty{21}{\micro m}، وحجم أربعين ضعف نواة الإنسان.
\textbf{23.} $\num{3.2e7}$ ثمانية، أي $\num{3.2e10}$ بقية --- أي عُشر
تصنيع \omterm{def:b1:proteins:peptide}{البروتين} اليومي في \omterm{def:b1:cell-unit-of-life:cell}{الخلية}، منفَقًا على التغليف.
\textbf{24.} \omterm{def:b1:genomes:genome}{صبغي} البكتيريا أقصر بألف مرة ويكفي \omterm{prop:b1:genomes:supercoiling}{الالتفاف الفائق} في
عرى لإدخاله في ميكرومتر؛ أما \omterm{def:b1:cell-unit-of-life:prokeuk}{حقيقية النواة} فعليها أن تطوي طولًا يفوق
نواتها بعشرة آلاف مرة وتحتاج إلى هرمية رص، \omterm{def:b1:genomes:nucleosome}{الجسيم النووي} أول مستوياتها.
\textbf{25.} نحو \num{8500} (أي \qty{8.5}{cm} في \qty{10}{\micro m})،
مبنيًا على $\times 6$ (الجسيمات النووية)، ثم $\times 6$ أخرى (الليف
\qty{30}{nm}، أي $\times 37$ في المجموع)، ثم $\times 27$ أخرى (العرى،
أي $\times 1000$ في المجموع)، ثم $\times 8$ أخرى (تكثف الطور
الاستوائي).
\end{solution}
